\section{R2. Carriers, accumulators and auxiliary state}\label{sec:carriers}

\subsection{The problem is a representation together with its operations}
\label{sec:carrier-package}

A Church list exposes its fold,
\code{xs : (R : Type) -> (A -> R -> R) -> R -> R}; using it requires choosing
$R$. Ordinary recursors expose the same choice through their motive. When the
requested result is already the fold result, unification may infer that type.
Otherwise the useful representation need not occur in the query. This motivates
carrier invention, but does not establish that every bounded search failure is
caused by a missing carrier: transition synthesis, finishing, proof and search
allocation must be measured separately.

For a list function $h:\mathrm{List}\,A\to T$, search a \emph{package}
\begin{equation}\label{eq:carrier-package}
 (K,z,s,o),\qquad z:K,\quad s:A\to K\to K,\quad o:K\to T,
 \qquad h(xs)=o(q(xs)),
\end{equation}
where $q([])=z$ and $q(a::xs)=s(a,q(xs))$. A direct fold has $K=T$ and
$o=\mathrm{id}$; a fold with a general finisher is a different problem.
The package may depend on fixed query parameters, but must not capture the
major input $xs$ outside its fold. Otherwise $K=\mathrm{Unit}$ can ignore the
fold while its finisher computes the answer from the captured input. Do not
allow calls to the unfinished target except through approved recursion
capabilities (N4's scope restriction).

Even without that capture, $K=\mathrm{List}\,A$, $z=[]$, $s=\mathrm{cons}$
and $o=h$ always reconstruct the input and defer the entire problem to the
finisher. Charge and bound \emph{all} package components, required helpers and
proofs. State cardinality, carrier-type syntax, executable description length,
proof cost, runtime and storage are different objectives. A natural number
can encode a pair of natural numbers; a short function type may contain
closures retaining input-sized data. Neither makes the corresponding encoding
cheap or useful.

\begin{center}\small
\begin{tabular}{@{}L{35mm}L{51mm}L{60mm}@{}}
\toprule
Task or scheme & Useful representation & Information or operation supplied\\
\midrule
\code{map}, \code{append}, \code{filter} & result type & direct fold\\
right-fold implementation of \code{foldl} & $T\to T$ & delayed accumulator\\
right-associated nonempty reduction & $\mathrm{Option}\,A$ & empty/nonempty phase\\
left-associated nonempty reduction & $\mathrm{Option}\,A\to A$ & phase and continuation\\
\code{maximum} with a default & $\mathrm{Option}\,A$ & seen element and comparison\\
\code{minMax} & $\mathrm{Option}\,(A\times A)$ & two extrema and empty phase\\
index with supplied default & $\mathrm{Int}\to A$ & consume an index through the fold\\
\code{fib} & $\Nat\times\Nat$ & adjacent recurrence values\\
efficient reversal/tree flattening & $\mathrm{List}\,A\to\mathrm{List}\,A$ & continuation/difference list\\
\bottomrule
\end{tabular}
\end{center}

These are templates, not semantic necessities or new measured successes.
In particular, do not assume that an optional element alone implements an
arbitrary nonassociative \emph{left}-associated operation through a right fold.
The historical targeted $T\to T$ rule and the cost of the blind carrier
frontier motivate keeping useful specializations early in a general grammar;
uniform enumeration need not erase them. The same problem appears as motive
generalization, tupling, invariant strengthening and helper invention: each
proposes information whose specification is not explicitly supplied.

\subsection{Minimal semantic state and what samples cannot identify}
\label{sec:carrier-quotient}

The classical direct-fold characterization requires that the kernel of $h$
respect \code{cons}: equal results on tails remain equal after prepending an
element~\cite{whenfold}. For a fold with hidden state and a finisher, the more
appropriate relation is
\begin{equation}\label{eq:carrier-residual}
 x\sim_h y\quad\Longleftrightarrow\quad
 \forall u:\mathrm{List}\,A,\quad
 h(u\mathbin{+\!+}x)=h(u\mathbin{+\!+}y).
\end{equation}
The contexts \emph{prepend} a list because a right fold has already processed
the tail. Reversing this orientation changes the state problem.

\begin{theorem}[Residual quotient and reachable state]
\label{thm:carrier-quotient}
For every realization~\eqref{eq:carrier-package}, $q(x)=q(y)$ implies
$x\sim_h y$. The quotient by $\sim_h$ supports a realization of $h$, and
the reachable states of any realization map surjectively onto that quotient.
In particular, $r$ pairwise inequivalent tails require $r$ distinct states.
\end{theorem}
\begin{proof}
Equal states remain equal after applying the same sequence of prepend
transitions; induction on $u$ gives equal states on $u\mathbin{+\!+}x$ and
$u\mathbin{+\!+}y$, and the common finisher gives equal outputs. For the
quotient, choose $z=[[]]$, $s(a,[x])=[a::x]$ and $o([x])=h(x)$.
The step is well defined because a test with prefix $u$ on $a::x$ is a test
with prefix $u\mathbin{+\!+}[a]$ on $x$; the empty prefix makes the output
well defined. Induction gives $q(x)=[x]$. Sending a reachable state $q(x)$
to $[x]$ is well defined by the first implication and reaches every class.
\end{proof}

This is a semantic characterization, not an effective minimal-type algorithm.
The quotient can be infinite, its equality undecidable, and its useful
representation absent from the chosen grammar. The finite lower bounds use
only finite counting; stronger comparisons of arbitrary infinite cardinalities
need their set-theoretic assumptions stated separately. Isomorphic
representations need not have the same syntax or operational cost.

Finite samples do not determine this universally quantified relation in
general. Over a nonempty alphabet, for any finite collection of input words
choose an unobserved word $w$. Constant zero
and the function returning one precisely on $w$ agree on zero-valued samples,
yet only the first has a one-state deterministic realization. Both functions
are total and computable. Thus samples do not determine even whether the
intended function requires more than one state.

\subsection{Two constructive bounded problems}
\label{sec:carrier-finite}

\begin{theorem}[Minimum finite machine consistent with finite samples]
\label{thm:carrier-finite-machine}
Let $A,B$ be finite explicitly enumerable sets with decidable equality and
$B$ nonempty. A finite consistent sample set
$D\subseteq\mathrm{List}\,A\times B$ admits a terminating procedure that
finds a minimum-state deterministic right-fold machine fitting $D$.
\end{theorem}
\begin{proof}
Let $S$ contain $[]$ and all suffixes of sample inputs. Use states
$S\cup\{\bot\}$, start at $[]$, and set $s(a,x)=a::x$ if $a::x\in S$,
otherwise $\bot$; the sink stays at $\bot$. Processing a sample from right
to left stays in $S$. Assign its specified output to each sample state and a
fixed $b_0\in B$ elsewhere. Consistency makes this well defined and gives
an upper bound $|S|+1$ on state count. Equivalently, construct a trie of
reversed sample words and add a sink.

For $k=1,\ldots,|S|+1$, enumerate every transition table
$A\times\mathrm{Fin}\,k\to\mathrm{Fin}\,k$ and output table
$\mathrm{Fin}\,k\to B$, fixing the seed to zero. Any nonempty finite
machine can rename its initial state to zero. Every sample check terminates;
the first successful $k$ is minimal and the upper bound ensures a success.
\end{proof}

With $a=|A|$ and $b=|B|$, there are $k^{ak}b^k$ tables at fixed seed, or
$k^{ak+1}b^k$ with all labeled seeds. These different conventions explain
otherwise different incoming enumeration counts. The theorem establishes
decidability, not interactive tractability. Direct backtracking with early
sample propagation is a Lean-only option; SAT or finite-domain solving,
symmetry breaking and partition refinement are alternative accelerators.
One encoding uses states $q_x$, transitions $d_{a,k}$ and outputs $o_k$ with
\[
 q_{a::x}=d_{a,q_x},\qquad o_{q_x}=h(x)
\]
where observations exist, preserving consistency whenever states are shared.
A satisfying finite table still needs a permitted typed emitter, finite
enumeration/equality operations and native checking. It is not automatically
expressible within an arbitrary small step grammar.

\paragraph{Bounded typed packages.}
Fix finite type and term productions, providers, literals, universes, helpers,
recursion schemes and proof choices, together with a total syntax grade for
$(K,z,s,o)$. Enumerate in a fixed objective order, type-check and validate every
observation, and prune only by valid necessary conditions. If the resulting
space $\mathcal P_b$ is finite and each check terminates decisively, this
returns a minimum sample-consistent member or proves none exists in
$\mathcal P_b$: every member is visited, no satisfying member is pruned, and
the first success is least in the chosen order. These are all needed premises.
For arbitrary Lean proof or reduction attempts, retain unknown outcomes;
``enumeration exhausted with unresolved checks'' is not semantic absence
(Section~\ref{sec:semantic-boundaries}). Universal contracts still require
their actual proofs.

Bounding only type syntax does not give this procedure. For prescribed
$K=X$ and target $\mathrm{Unit}$, even the single empty-list observation has
a realization exactly when $X$ has a closed inhabitant: the seed supplies one,
and conversely a chosen seed can be preserved by the step and discarded by
the finisher. Conversely an unbounded finisher $\mathrm{Unit}\to T$ already
contains arbitrary inhabitation of $T$. These reductions from N8 and N9 show
why type enumeration alone cannot decide arbitrary package existence.

\subsection{Sound sample lower bounds and their counterexamples}
\label{sec:carrier-lower-bounds}

Build a witnessed incompatibility graph on selected tails. Join $x$ and $y$
only when a common prefix $u$ has two known, different outputs for
$u\mathbin{+\!+}x$ and $u\mathbin{+\!+}y$.

\begin{proposition}[Observed-context capacity bound]
\label{prop:carrier-capacity}
Any deterministic package fitting these observations assigns different states
to adjacent tails. For finite $K$,
\[
 \chi(G)\le |\operatorname{Reach}(K)|\le |K|.
\]
A clique of size $r$ therefore refutes a carrier with fewer than $r$ possible
reachable states.
\end{proposition}
\begin{proof}
If $q(x)=q(y)$, processing the common prefix gives equal states and equal
finished outputs, contradicting its witnessed unequal observations. Hence
$x\mapsto q(x)$ is a proper coloring by reachable states.
\end{proof}

The certificate consists of the actual conflicting observations and a proved
capacity bound, or an inductively closed invariant bounding the reachable
state set for \emph{every} admitted program in the claimed class. A clique is
cheap to check. A general chromatic lower bound needs its own evidence;
the number of colors found by greedy coloring is an \emph{upper} bound on
$\chi(G)$, not a refutation certificate. Even a correct coloring need not
extend to a deterministic transition system: equal states must have equal
successor states under the same constructor.

Three further mistakes recur in the incoming answers' counterexamples.
\begin{itemize}
\item Partial rows $(0,?)$, $(?,1)$, $(1,1)$ show that compatibility is not
transitive. Unknown entries are neither outputs nor inequalities. Different
partial rows may describe the same state; merging all compatible pairs or
counting distinct strings of entries is unjustified.
\item Values found by a depth-limited enumerator are not an upper bound on
inhabitants or reachable states. A short transition can build arbitrarily
large values when iterated. Refuting \code{Nat}, lists or functions because
only two values were found is unsound without a separate bounded-image proof.
\item A direct-fold obstruction does not exclude the same carrier with a
finisher. For total list tail, $h([])=h([b])=[]$, but after prepending $a$,
$h([a])=[]\ne[b]=h([a,b])$. This excludes a direct fold whose state is
$h(xs)$. The same list carrier with $q=\mathrm{id}$ and $o=\mathrm{tail}$
works, so the restriction on finishing code is essential.
\end{itemize}

\paragraph{Three states for a Boolean result.}
For ``length divisible by three'', tails of lengths $0,1,2$ under prefix
lengths $0,1,2$ have rows
$(1,0,0)$, $(0,0,1)$, $(0,1,0)$. They form a three-clique. A cyclic
three-state counter with output true at zero attains the bound; induction
identifies the state with length modulo three. Thus \code{Bool} is too small
even though the output is Boolean; \code{Option Bool} supplies three values.
Unary fixed-seed table enumeration uses $2,16,216$ tables at sizes one, two
and three; changing the alphabet or allowing every labeled seed changes these
counts. N1, N3, N4 and N9 contain variants of this finite illustration.

\paragraph{Four states from nine distinct observations.}
Let $h$ mean that a word over $\{a,b\}$ contains both symbols. The tails
$[],[a],[b],[a,b]$ under prefixes $[],[a],[b]$ yield
\[
\begin{array}{c|ccc}
\text{tail} & [] &[a]&[b]\\\hline
{[]}&0&0&0\\
{[a]}&0&0&1\\
{[b]}&0&1&0\\
{[a,b]}&1&1&1
\end{array}
\]
Every pair is incompatible. The twelve cells use nine distinct input words;
the stronger sample set of all fifteen words of length at most three contains
them. Four states suffice: retain two seen-symbol flags, set the relevant flag
on each step and finish with conjunction. An induction establishes the flag
invariant on all words. Fixed-seed enumerations of the smaller binary machines
have $2,64,5832$ tables; including all seeds gives $2,128,17496$.
N2, N5 and N7 use different sample extensions and witness-test bounds. These
are compatible finite experiments, not interchangeable receipts. Their Python
checks corroborate the construction, not native Lean synthesis, universal
correctness or performance; the mathematical induction supplies the universal
argument. N4's binary parity example similarly needs two states, demonstrating
that these bounds concern the particular observed behavior.

\paragraph{A minimum sample machine can still overfit.}
\label{sec:carrier-minimum-overfit}
N8's evaluation section supplies a distinct constructive example. For the unary
target ``length exactly two'', initially observe lengths zero through five.
Suffix lengths $0,1,2,3$ form a four-clique: each pair has an observed common
prefix taking exactly one member to length two. The four-state cycle
$0\to1\to2\to3\to0$, with output true only at state two, attains this
sample minimum but returns true incorrectly at length six. After adding the
length-six counterexample, the machine
$0\to1\to2\to3\to3$ with the same output convention still uses four states
and fits the strengthened samples. Its reached state is $\min(n,3)$ by
induction, proving the intended behavior at every length. Thus certified state
capacity, minimum sample consistency and identification of the intended
function are three different claims even in a tiny exact finite language.
The extra counterexample comes from the experiment's reference oracle, which
an arbitrary example-only query does not automatically provide.

Active automata learning supplies useful observation-table ideas, but Angluin's
membership and equivalence oracles are stronger than a fixed sample set~\cite{angluin}.
Do not assume missing table entries, closure or consistency have been learned.
Use state constraints to suggest options, products or continuations; the table
is not a decoder of an intended Lean type.

\subsection{Polymorphic carriers: shapes, positions and available observations}
\label{sec:carrier-parametricity}

The correct reduction is conditional on a relational or naturality theorem,
not merely a polymorphic-looking Lean type. The incoming reports connect
polymorphic testing~\cite{bernardy10,wadler89} with Mulleners, Jeuring and
Heeren's container-based realizability work~\cite{containers}. N9 supplies the general
construction; the list instance used by the other answers follows from it.

\begin{theorem}[Shape--position representation]
\label{thm:carrier-containers}
For containers $F(A)=\sum_{s\in S}(P_s\to A)$ and
$G(A)=\sum_{t\in U}(Q_t\to A)$, a natural transformation
$\tau:F\Rightarrow G$ is determined by a shape map $\phi:S\to U$ and
maps $\psi_s:Q_{\phi(s)}\to P_s$, with
\[
 \tau_A(s,k)=(\phi(s),k\circ\psi_s).
\]
Conversely these maps define a natural transformation.
\end{theorem}
\begin{proof}
Evaluate $\tau$ at the generic input $(s,\mathrm{id}_{P_s})$ of element
type $P_s$; its result defines $(\phi(s),\psi_s)$. Naturality along
$k:P_s\to A$ gives the displayed equation. Conversely, associativity of
composition makes this equation commute with every element map.
\end{proof}

For lists the shape is a length $n$ and positions are $\mathrm{Fin}\,n$.
If $f_A:\mathrm{List}\,A\to\mathrm{List}\,A$ is natural, let
$e_n=[0,\ldots,n-1]$ and let $v:\mathrm{Fin}\,n\to A$ select an input
list's entries. Then
\begin{equation}\label{eq:carrier-position-test}
 f_A(xs)=\operatorname{map}(v)(f_{\mathrm{Fin}\,n}(e_n)).
\end{equation}
The output word records deletion, duplication or permutation of positions.
Repeated elements are covered because $v$ need not be injective. At $n=0$
the position type is empty and its list result is empty; no nonexistent
function $\Nat\to A$ or arbitrary element of an empty type is assumed. Use
appropriate universe lifts at a higher fixed universe. For multiple containers
tag position sets by argument and retain specified cross-input sharing.

If two natural functions agree on $e_n$ for $n\le N$, applying
\eqref{eq:carrier-position-test} proves equality on every input of length at
most $N$. This does not prove equality at greater lengths or for all programs
of a syntax-size bound without a separate small-model theorem. Identity and
the natural function that returns its input through length $N$ and then
returns $[]$ agree on those tests and disagree at $N+1$.

Nor does quantifying over \code{Type} establish naturality for arbitrary Lean
providers. Classical decidability permits a noncomputable family returning
its input when $A$ is subsingleton and $[]$ otherwise. Mapping a singleton
list between one- and two-element types witnesses failure of naturality.
An axiom-profile check alone is not a parametricity certificate.

With lawful equality, arbitrary maps that collapse elements need not preserve
behavior; bounded inputs require realizable partitions of positions, not just
the all-distinct case. Orders, hashing, algebraic dictionaries, distinguished
elements and callbacks expose further structure. A raw \code{BEq} dictionary
need not be lawful; retain its operational semantics until a preservation law
permits replacing it with mathematical equality. Nominal or register-style
models can separate finite control from stored atoms and symbolic payloads;
nominal automata provide a relevant semantic precedent~\cite{nominal},
but reversing or copying arbitrarily long lists may require unbounded storage.
Finite symbolic alphabets do not imply finite concrete state.

A practical provider classification is token-preserving, shape-observing,
equality-observing, order-observing or opaque. Certify the admitted fragment
and providers before using generic positions for permanent pruning or merging.
Without that theorem, symbolic observations remain useful tests and proposal
features. Indexed tokens additionally carry their index and typing telescope;
renaming must preserve fibers, and caching modulo token renaming needs an
evaluator invariance theorem.

\subsection{Constructive fusion, feature discovery and closure}
\label{sec:carrier-lifting}

The original claim that auxiliary-state synthesis has no relevant published
algorithm is too broad. AutoLifter synthesizes auxiliary summaries and
combinators from a reference computation and a prescribed algorithmic paradigm;
its decomposition can choose necessary conditions with unrealizable later
subtasks. It is neither a global minimal-type method nor a complete sparse-
example learner~\cite{autolifter}. The incoming reports also identify Eguchi and colleagues'
auxiliary-type inference, Para's recursion templates, and Synduce's
unrealizability witnesses as useful but differently scoped precedents~\cite{eguchi,para,synduce}.
Tupling and homomorphism calculation remain relevant~\cite{tupling,thirdhom}.

The earlier comparisons remain useful when their task boundaries are retained.
$\lambda$\textsuperscript{2} propagates examples into a fixed combinator vocabulary;
its fold deduction uses tail-complete observations and a result-typed
accumulator~\cite{lambda2}. Myth, Smyth, Burst and Trio investigate explicit
recursion, trace information, live evaluation, angelic results or construction
from nonrecursive blocks~\cite{myth,smyth,burst,trio}. Those mechanisms do not
by themselves constitute the joint carrier/package learner specified here.
Synquid's supplied auxiliary refinements and DreamCoder's abstractions learned
across a corpus address different sources of missing specification
\cite{synquid,dreamcoder}. These distinctions identify adaptation opportunities;
they do not support the original universal claim that synthesis research never
invents auxiliary state.

\paragraph{An evidence-dependent constructive fold.}
Weber and Caldwell's constructive work qualifies the idea that fold
characterization is purely nonconstructive~\cite{weber}. The following sufficient condition
from N6 makes the supplied evidence explicit.

\begin{proposition}[Fold construction from a section]
\label{prop:carrier-section-fold}
Let $h:\mathrm{List}\,A\to B$ have a section
$r:B\to\mathrm{List}\,A$ satisfying $h(r(b))=b$. If
$h(x)=h(y)$ implies $h(a::x)=h(a::y)$, then
\[
 h=\operatorname{foldr}\ s\ z,\qquad
 z=h([]),\qquad s(a,b)=h(a::r(b)).
\]
\end{proposition}
\begin{proof}
The base case is immediate. Since $h(r(h(x)))=h(x)$, cons congruence gives
$h(a::r(h(x)))=h(a::x)$. Substitute the induction hypothesis for the fold
of the tail to obtain the constructor equation.
\end{proof}

This consumes $h$, a section and a congruence proof; it does not discover them
from samples. A section on the reachable image, with its typing and closure
obligations explicit, can weaken the unnecessarily global surjectivity premise.
For a polynomial functor the corresponding construction applies $F(r)$ before
the constructor and $h$, with the analogous congruence and induction premises.

\paragraph{Feature separation is only the first stage.}
Choose a bounded typed feature library $\phi_j:\mathrm{List}\,A\to S_j$:
emptiness, optional first/last element, length modulo a small modulus, flags
for a concrete alphabet, or available measures. Features must use only the
actual operations of the query; abstract elements cannot be compared with an
unsupplied distinguished value. Each feature separates some observed
incompatibility edges. Weighted set cover proposes a small set $J$ and
\[
 K_J=\prod_{j\in J}S_j,\qquad
 \phi_J(xs)=(\phi_j(xs))_{j\in J}.
\]
Now synthesize the seed, updater and finisher with the closure equations
\begin{align}
 \phi_J([])&=z,&
 \phi_J(a::xs)&=s(a,\phi_J(xs)),&
 o(\phi_J(xs))&=h(xs).\label{eq:carrier-feature-closure}
\end{align}
Induction on $xs$ proves $\operatorname{foldr}\ s\ z\ xs=\phi_J(xs)$
from the first two equations; the third gives the requested result. For a
partially specified task retain the actual contract instead of inventing a
reference equality. Finite observations may suggest the equations but do not
prove their universal versions.

If two tails agree on selected features but need different successor features
after the same constructor, add a closure-distinction obligation and seek more
information. A feature set that separates outputs may still be impossible to
update cheaply. Retain alternative sets and backtrack; greedy feature selection
is a heuristic, not a complete lifting procedure. Exhaustively enumerate
bounded feature subsets, algebras and finishers with complete checking for a
relative completeness theorem. Preserve unknown proof alternatives. This is
the constructive separation-and-closure procedure emphasized by N2, N3 and N7.

Alternatively propose a relation $I(xs,k)$ describing stored information and
prove initialization, preservation and finishing:
\[
 I([],z),\quad
 I(xs,k)\Rightarrow I(a::xs,s(a,k)),\quad
 I(xs,k)\Rightarrow C(xs,o(k)).
\]
The invariant is itself a synthesis proposal. A reference computation may
generate counterexamples; sparse contracts provide only their stated
observations. For $F(X)=\sum_c A_c\times X^{r_c}$, each constructor receives
an algebra $A_c\times K^{r_c}\to K$, using the summaries of \emph{all}
recursive children. The general closure equation is
$H\circ\mathrm{in}=\alpha\circ F(H)$; for a lifted
$H=\langle h,b\rangle$, projecting the fold gives $h$. A fixed carrier and
bounded constructor algebra give a finite synthesis class. Polynomial functors
alone do not make unrestricted helper and invariant invention decidable.

\subsection{Worked representations and scheme-dependent invariants}
\label{sec:carrier-worked}

\paragraph{Reversal is already a list-valued fold.}
\label{sec:carrier-reverse}
All nine answers correct the semantic use of reversal in the original proposal:
\[
 \operatorname{reverse}(xs)=
 \operatorname{foldr}\ (\lambda a\ r.\ r\mathbin{+\!+}[a])\ []\ xs.
\]
The induction step is exactly the defining reversal equation. A shallow or
append-free grammar may not express the step, but that is a language
restriction, not a capacity obstruction. Under ordinary singly linked lists,
the successive append traversals cost $0+1+\cdots+(n-1)=\Theta(n^2)$.
For the continuation representation set
\[
 K=\mathrm{List}\,A\to\mathrm{List}\,A,\quad z=\mathrm{id},\quad
 s(a,k)=\lambda acc.\ k(a::acc),\quad o(k)=k([]).
\]
Induction proves $q(xs)(acc)=\operatorname{reverse}(xs)\mathbin{+\!+}acc$:
the step uses the induction hypothesis at $a::acc$ and append associativity.
With constant-time cons and closure creation, construction and final traversal
of the closure chain take linear work. A left-fold scheme can also reverse
linearly with an ordinary list carrier. These alternatives show why the
scheme, operation grammar and cost model belong in the optimization problem.

\paragraph{Adjacent values, empty phases and defaults.}
For Fibonacci, $q(n)=(F_n,F_{n+1})$, seed $(0,1)$ and transition
$(a,b)\mapsto(b,a+b)$ close the recurrence; induction proves the invariant.
This does not prove that no clever single-natural encoding exists. Optional
maximum state avoids inventing an element before one has been seen, and an
optional pair similarly handles extrema. A total maximum, nonempty reduction
or index must preserve the supplied empty/out-of-range default or inhabitance
assumption. Carrier invention cannot create an element of an uninhabited
type. Difference lists for tree flattening avoid repeated copying under the
same explicit operational assumptions used for reversal.

\paragraph{Maximum prefix sum: the recursion scheme changes the needed state.}
\label{sec:carrier-prefix-sum}
N9 supplies a useful distinction not captured by the earlier examples. For
integer lists, let $h$ be maximum prefix sum including the empty prefix. A
right fold needs only an integer:
\[
 h([])=0,\qquad h(a::xs)=\max(0,a+h(xs)).
\]
But the result alone cannot combine arbitrary adjacent segments: $[0]$ and
$[-1]$ both give zero, while appending $[1]$ gives one and zero respectively.
Lift to $H(xs)=(\operatorname{sum}(xs),h(xs))$, with
\begin{equation}\label{eq:carrier-prefix-product}
 (s,m)\otimes(t,n)=(s+t,\max(m,s+n)).
\end{equation}
Every prefix of $x\mathbin{+\!+}y$ lies inside $x$ or contains all of $x$
and a prefix of $y$, so $H(x\mathbin{+\!+}y)=H(x)\otimes H(y)$.
The reachable invariant is $I(s,m):\ 0\le m\land s\le m$.

\begin{proposition}[Invariant-restricted summary monoid]
\label{prop:carrier-prefix-monoid}
Operation~\eqref{eq:carrier-prefix-product} is associative on $\mathbb Z^2$,
preserves $I$, and has two-sided identity $(0,0)$ on states satisfying $I$.
\end{proposition}
\begin{proof}
Both bracketings have first coordinate $s+t+u$ and second coordinate
$\max(m,s+n,s+t+p)$. For invariant operands the second coordinate is at least
$m\ge0$ and at least $s+t$ because $n\ge t$. Left identity uses
$\max(0,m)=m$; right identity uses $\max(m,s)=m$.
\end{proof}

Left identity fails on the unrestricted pair $(1,-1)$. Demanding algebra laws
for all syntactic pairs can therefore reject this correct reachable-state
representation; assuming $I$ without proving its initialization and preservation
can validate an incorrect one. Carrier type, scheme and invariant are related
search choices. ``Minimal carrier'' has no scheme-independent meaning here.

\subsection{Motive enumeration and a shared strengthening service}
\label{sec:carrier-motives}

Retain a short prioritized carrier grammar, for example
\[
 K::=T\mid A\mid S_i\to K\mid K\times K\mid\mathrm{Option}\,K
       \mid\mathrm{Bool}\mid\Nat,
\]
then its indexed and dependent extensions when needed. A size-three frontier
is an engineering starting point, not a theorem of coverage. Every proposal
opens seed, algebra and finishing obligations together; a type that elaborates
but cannot be initialized is not a solved representation. Charge the whole
subtree and preserve specialized priorities such as $T\to T$.

For $f:\Nat\to S_1\to\cdots\to S_k\to T$, a constant recursor result
$K=S_1\to\cdots\to S_k\to T$ lets the successor branch apply its
predecessor result at newly computed parameter values. This is a complete
schema for the declared nondependent primitive-recursive fragment provided
its entire base and step language is enumerated. It is not a small complete
procedure for every useful dependent primitive recursion.

Separate existing-parameter generalization, result strengthening, index
abstraction and recursion nesting. A dependent template is
\[
 M(\vec\imath,x)=
 \prod_{j\in J}(a_j:S_j(\vec\imath,x,\vec a_{<j})).\,
 K(\vec\imath,x,\vec a),
\]
with $J$ dependency-closed and binders in valid telescope order. Dependent
accumulators need $\Pi$ rather than ordinary arrows; coupled values and
proof-carrying summaries may need $\Sigma$ rather than plain products.
Type-check motives before opening all method searches, and deduplicate only
with justified equality, not a display-level index simplification.

Dovetail over generalized binders, carrier size, total term size, helper
count, active recursion frames and proof/universe bounds. One possible joint
grade charges $|K|+|M|+|z|+|s|+|o|$ plus recursive-scheme introductions.
Fair expansion covers the union of the declared finite grammars under the
appropriate checking assumptions, not all Lean inhabitants at an interactive
deadline. Keep outermost natural-number recursion and existing arithmetic
providers early, use observed constants and a small literal basis first, and
charge nested eliminators separately. Permanently excluding later schemes or
numerals loses the corresponding completeness claim. A beam or a fixed top-$k$
carrier list needs a fair fallback retaining its omissions.

\paragraph{Failure-directed strengthening.}
Carrier choice, motive generalization and helper invention should share a
bounded service that consumes structured failed obligations. Record the
available induction hypothesis, required recursive application, changed
parameters, unmatched indices, blocked equations and distinguishing residual
observations. Generalize a changing accumulator and its dependencies; add a
missing statistic as a product; add an empty phase tag; propose a typed helper
for a repeated fragment; or retrieve an algebraic reassociation lemma when
that is the actual missing step. A helper that merely adds another unsolved
function is not useful until both output recovery and closure become easier.

Use typed anti-unification and rippling-inspired comparison to generate a
small ranked menu, preserving the original branch. Tests reject false
conjectures cheaply; survivors remain conjectures until native proof and
policy checking. HipSpec and Hipster supply theory-exploration precedents~\cite{hipspec,hipster},
not a turnkey Lean import established by these reports. Induction with
generalization in superposition reasoning is another supplied precedent for
strengthening induction problems~\cite{inductiongeneralization}; transferring
its first-order inference machinery requires a separate typed Lean design.
Canonical's supplied
generalized induction example demonstrates predicate discovery with a suitable
principle available~\cite{canonical}; discovering that principle itself is a
separate task. If a guessed helper fails, backtrack its consumers. Record
environment and axiom profile before reusing a proved lemma across tasks.

\subsection{Evaluation and new expert questions}
\label{sec:carrier-experiments}

Use at least three matched tracks: carrier supplied, carrier selected from a
fixed template grammar, and carrier inferred from failures or observations.
Add an efficiency track when runtime/storage, rather than merely sample
consistency, is the objective. Hold the original contract, defaults, allowed
providers, universe policy and body grammar fixed across relevant comparisons.
Compare blind and distinction-guided ordering of the same representation
grammar; measure subtree size and time per carrier, transition/finisher search,
proof debt and final code cost. Scheme selection, motive discovery, body
construction, universal proof and executable publication are separate results.

The historical thirteen Church stretch cases and accumulator tasks such as
\code{fib}, \code{list_rev_tailcall}, \code{tree_flatten}, \code{maximum}
and \code{index} remain useful evaluation families. The original goal of
solving them within ten seconds without new task-specific rules and without
more than a factor-two slowdown of solved Church probes is a prospective
acceptance target, not a claim supplied by the incoming finite models. No
exact modern accuracy or speedup follows from their Python checks.

The historical five expert questions now have the conditional answers and
explicit limits above. The active next questions are:
\begin{description}
\item[R2.N1.] Can a Lean-native carrier proposer emit and cheaply check
distinguishing-context certificates against finite capacity or a proved
reachable-state invariant, while sharing work with bounded algebra synthesis?
\item[R2.N2.] Which provider-certified parametric fragment supports a proved
shape/position or equality-pattern reduction for higher-order and indexed
queries, and can a useful small-model bound connect program size to test depth?
\item[R2.N3.] Can failure-directed feature separation and closure synthesis
retain relative completeness while materially reducing search over the same
carrier, updater and finisher grammar on sparse-contract tasks?
\item[R2.N4.] What joint grammar and cost model can select recursion scheme,
dependent motive, carrier and reachable-state invariant without pushing the
original computation into an unrestricted finisher or helper?
\item[R2.N5.] Which structured induction failures reliably distinguish a missing
parameter, missing state component and missing algebraic lemma, and how much
does that diagnosis improve discovery after charging its extra proof debt?
\end{description}
